\chapter{Intro}
\section{Topological manifolds}
\lecture{1}{Wed 03 Sep 2025 12:16}{Syllabus day}
Smooth manifolds form a proper subset of topological manifolds. Interesting.

\begin{definition}\label{1.1.1}
	A topological manifold is a Hausdorff topological space $(M,\mathcal{T} )$ which is second countable (it has a countable basis), and for some fixed $n\in\mathbb{N} $, for all points $p\in M$ there is an open set $p\in U\subseteq M$ and a homeomorphism $f : U \cong \hat{U} \subseteq \mathbb{R} ^n$ for some $\hat{U} $. We say the dimension of $M$ is $n$, and call it an $n$-manifold.
\end{definition}

Of course we may equivalently say for all points $p$, there is an open neighborhood $U$ and a homeomorphism $f : U \cong \mathbb{R} ^n$. This follows since $\hat{U} $ admits an open cover by open balls, who have open preimages, and thus at least one has an open preimage $V$ containing $p$. Since the image of $f|V$ is an open ball, and open balls are homeomorphic to $\mathbb{R} ^n$, the conclusion follows.
